\editionPage[\pdfbookmark[0]{Corrigés des bilans}{bm-bilan-A-cor-auto-parent}]{bilan-A-cor-auto}{Bilan A · Corrigé}{Automatismes expliqués}{\hyperref[bilan-A-auto]{Énoncé p.\,\pageref*{bilan-A-auto}}}
\begin{blocEdition}{bilan-A-cor-auto}
\begin{questions}
\item Mêmes signes, produit positif : $\res{24}$.
\item Le produit est prioritaire : $7-10=\res{-3}$.
\item $\frac68+\frac18=\res{\frac78}$.
\item L’inverse est $\res{-\frac52}$ : le produit des deux nombres vaut 1.
\item $3+4=7$ facteurs 2 : $\res{2^7}$.
\item $11^2=121$ et $11$ est positif, donc $\res{11}$.
\item On distribue 3 : $\res{6x-15}$.
\item $4x=16$, donc $\res{x=4}$.
\item $\frac{15}{100}\times80=\res{12}$.
\item Le coefficient multiplicateur vaut $1+\frac8{100}=\res{1{,}08}$.
\item La troisième valeur est centrale : $\res8$.
\item $1-0{,}7=\res{0{,}3}$.\end{questions}
\end{blocEdition}
\editionPage{bilan-A-cor-1}{Bilan A · Corrigé 1}{Deux ateliers}{\hyperref[bilan-A-ex-1]{Énoncé p.\,\pageref*{bilan-A-ex-1}}}
\begin{blocEdition}{bilan-A-cor-1}
\editionEtape{1 · Modéliser}$A(n)=12+4n$ et $B(n)=7n$, en euros. Le premier tarif n’est pas proportionnel au nombre de séances, à cause de la partie fixe.
\editionEtape{2 · Résoudre}$12+4n=7n$ donne $12=3n$, donc $\res{n=4}$. Les deux dépenses valent alors $\res{28}$ euros.
\editionEtape{3 · Appliquer la remise}$A(8)=12+32=44$ euros ; avec le coefficient 0,9, on paie $44\times0{,}9=\res{39{,}60}$ euros. $B(8)=56$ euros. Le premier atelier est moins cher de $\res{16{,}40}$ euros.
\end{blocEdition}
\editionPage{bilan-A-cor-2}{Bilan A · Corrigé 2}{Un terrain et un arroseur}{\hyperref[bilan-A-ex-2]{Énoncé p.\,\pageref*{bilan-A-ex-2}}}
\begin{blocEdition}{bilan-A-cor-2}
\editionEtape{1 · Pythagore}ABC est rectangle en A : $BC^2=6^2+8^2=100$, donc $\res{BC=10\ \mathrm m}$.
\editionEtape{2 · Le cercle circonscrit}Le point équidistant des trois sommets est le centre du cercle circonscrit. Comme ABC est rectangle en A, c’est le $\res{\text{milieu de }[BC]}$, et la portée vaut $\frac{10}2=\res{5\ \mathrm m}$.
\editionEtape{3 · Les réserves}La longueur achetée est $10+2\times0{,}5=11$ m. Le coût vaut $11\times3{,}20=\res{35{,}20}$ euros. Les réserves s’ajoutent après le calcul de la distance.\label{fig-bilan-A-cor-2}\figureGeo[corrige]{bilan-a-geometrie}

\end{blocEdition}
\editionPage{bilan-A-cor-3}{Bilan A · Corrigé 3}{Une petite enquête}{\hyperref[bilan-A-ex-3]{Énoncé p.\,\pageref*{bilan-A-ex-3}}}
\begin{blocEdition}{bilan-A-cor-3}
\editionEtape{1 · Deux indicateurs}La somme vaut 205 min. La moyenne vaut $\frac{205}{12}\approx\res{17{,}1\ \mathrm{min}}$. Les 6e et 7e valeurs sont 15 et 20 : toute valeur entre 15 et 20 est une médiane, par exemple $\res{17{,}5\ \mathrm{min}}$.
\editionEtape{2 · Dispersion et fréquence}L’étendue vaut $30-5=\res{25\ \mathrm{min}}$. Six élèves lisent au moins 20 min : fréquence $\frac6{12}=\res{\frac12}$, soit 50\,\%.
\editionEtape{3 · Le hasard}Les six autres élèves lisent moins de 20 min. Chaque élève ayant la même chance d’être choisi, la probabilité vaut $\res{\frac12}$. L’événement contraire, « au moins 20 min », a pour probabilité $\frac12$ et $1-\frac12=\frac12$.
\end{blocEdition}
\editionPage{bilan-A-cor-4}{Bilan A · Corrigé 4}{Un cône et un cylindre}{\hyperref[bilan-A-ex-4]{Énoncé p.\,\pageref*{bilan-A-ex-4}}}
\begin{blocEdition}{bilan-A-cor-4}
\editionEtape{1 · Volume}$V=\frac{\pi\times3^2\times12}3=\res{36\pi\ \mathrm{cm}^3}\approx\res{113{,}1\ \mathrm{cm}^3}$.
\editionEtape{2 · Comparaison}Le cylindre a pour volume $\pi\times9\times12=108\pi$ cm$^3$, soit trois fois le volume du cône : il faut $\res3$ cônes remplis.
\editionEtape{3 · Contenance}Un cm$^3$ vaut un mL : la capacité est d’environ 113,1 mL. Donc $\res{110\ \mathrm{mL}}$ peuvent y tenir, si les dimensions intérieures sont celles données.\label{fig-bilan-A-cor-4}\figureGeo[corrige]{bilan-a-cone}

\end{blocEdition}
\editionPage{bilan-B-cor-auto}{Bilan B · Corrigé}{Automatismes expliqués}{\hyperref[bilan-B-auto]{Énoncé p.\,\pageref*{bilan-B-auto}}}
\begin{blocEdition}{bilan-B-cor-auto}
\begin{questions}
\item Signes contraires : $\res{-5}$.
\item On simplifie les facteurs : $\res{\frac32}$.
\item On divise par 12 : $\res{\frac34}$.
\item $(5\times2)^3=\res{10^3}$.
\item $81<90<100$, donc $\res{9<\sqrt{90}<10}$.
\item $\res{5x+5}$.
\item Le facteur commun est 6 : $\res{6(2x+3)}$.
\item $3x=12$, donc $\res{x=4}$.
\item $0{,}025\times1\,000=\res{25\ \mathrm L}$.
\item 7 parts de 5 : $\res{10\text{ et }25}$.
\item Une issue sur quatre : $\res{\frac14}$.
\item $(4-2\,;1+3)=\res{(2\,;4)}$.\end{questions}
\end{blocEdition}
\editionPage{bilan-B-cor-1}{Bilan B · Corrigé 1}{Un bassin et deux unités}{\hyperref[bilan-B-ex-1]{Énoncé p.\,\pageref*{bilan-B-ex-1}}}
\begin{blocEdition}{bilan-B-cor-1}
\editionEtape{1 · Convertir}$0{,}004$ m$^3$/s $=\res4$ L/s, soit $\res{240}$ L/min.
\editionEtape{2 · Durée}$t=\frac Vq=\frac{720}{240}=\res3$ min. On suppose le débit constant jusqu’à la fin.
\editionEtape{3 · Volume restant}En 90 s, la pompe retire $4\times90=360$ L. Il reste $\res{360\ \mathrm L}$, soit $\frac{360}{720}=\res{\frac12}$ de la capacité initiale.
\end{blocEdition}
\editionPage{bilan-B-cor-2}{Bilan B · Corrigé 2}{Des milieux et un parallélogramme}{\hyperref[bilan-B-ex-2]{Énoncé p.\,\pageref*{bilan-B-ex-2}}}
\begin{blocEdition}{bilan-B-cor-2}
\editionEtape{1 · Droite des milieux}Dans ABC, I et J sont les milieux de $[AB]$ et $[AC]$ : $\res{(IJ)\parallel(BC)}$ et $IJ=\frac{12}2=\res{6\ \mathrm{cm}}$.
\editionEtape{2 · Parallélogramme}$(IJ)\parallel(BK)$ car K est sur $[BC]$, et $IJ=6=BK$ car K est le milieu de $[BC]$. Le quadrilatère non croisé IBKJ a deux côtés opposés parallèles et de même longueur : c’est un $\res{\text{parallélogramme}}$.
\editionEtape{3 · Périmètre}$IB=3{,}5$ cm et $BK=6$ cm : le périmètre vaut $2\times(3{,}5+6)=\res{19\ \mathrm{cm}}$.\label{fig-bilan-B-cor-2}\figureGeo[corrige]{bilan-b-milieux}

\end{blocEdition}
\editionPage{bilan-B-cor-3}{Bilan B · Corrigé 3}{Une série et un tirage}{\hyperref[bilan-B-ex-3]{Énoncé p.\,\pageref*{bilan-B-ex-3}}}
\begin{blocEdition}{bilan-B-cor-3}
\editionEtape{1 · Statistiques}La moyenne vaut $\frac{6+16+24}{20}=\res{2{,}3}$. Dans la série ordonnée, les 10e et 11e valeurs valent toutes deux 2 : la médiane est $\res2$.
\editionEtape{2 · Probabilités}Il y a $8+6=14$ cartes paires sur 20 : $P(\text{pair})=\res{\frac7{10}}$. L’événement contraire a pour probabilité $\res{\frac3{10}}$.
\editionEtape{3 · Modifier le total}Les 14 cartes paires ne changent pas. Pour représenter la moitié, elles doivent être dans un total de 28 cartes : on ajoute $28-20=\res8$ cartes portant 1. Vérification : $\frac{14}{28}=\frac12$.
\end{blocEdition}
\editionPage{bilan-B-cor-4}{Bilan B · Corrigé 4}{Tripler une variable}{\hyperref[bilan-B-ex-4]{Énoncé p.\,\pageref*{bilan-B-ex-4}}}
\begin{blocEdition}{bilan-B-cor-4}
\editionEtape{1 · Suivi}À partir de 2, les tours donnent $\res{6,\ 18,\ 54,\ 162}$. La valeur initiale 2 n’est pas un tour.
\editionEtape{2 · Puissance}Quatre multiplications par 3 donnent $2\times3^4=2\times81=\res{162}$.
\editionEtape{3 · Remonter le calcul}Avec une entrée $a$, l’affichage vaut $81a$. L’équation $81a=405$ donne $\res{a=5}$. Contrôle : $5\times3^4=405$.\label{fig-bilan-B-cor-4}\figureGeo[corrige]{bilan-b-script}

\end{blocEdition}
